\section{总结与延伸}

\subsection{一条从 U-Net 到下一代生成系统的主线}

本讲可压缩为七次结构重写：把 diffusion pipeline 与 denoiser 分开；把巨型 U-Net 统一成 token backbone；用 AdaLN-Zero 注入 timestep/class；用 cross-attention 处理文本；用 linear attention 应对高分辨率；用 MMDiT 让模态共同演化；再用参数共享、控制网络与分层表示把成本压回可部署范围。

\begin{importantbox}{核心结论}
DiT 的真正贡献不是“Transformer 也能画图”，而是把生成 backbone 接入可扩展的 Transformer 设计空间。此后每个新模型都在回答同一组问题：token 怎样形成、条件怎样进入、模态是否共享参数、注意力如何缩放、以及质量收益是否值得系统成本。
\end{importantbox}

\subsection{工程检查表}

\begin{enumerate}
  \item \textbf{Pipeline：}是否同时报告 text encoder、VAE、scheduler 与 denoiser？
  \item \textbf{Conditioning：}低维条件用 AdaLN，长序列条件是否需要 cross/joint attention？
  \item \textbf{Resolution：}token 数多大，$N^2$ 何时成为瓶颈？
  \item \textbf{Sharing：}参数共享降低的是权重、FLOPs 还是实际 latency？
  \item \textbf{Evaluation：}是否有固定 prompt、失败样本、人类偏好与控制遵循指标？
  \item \textbf{Video：}是否测 temporal consistency，而非只看单帧？
  \item \textbf{Reproducibility：}数据、配置、权重、seed 和评估脚本是否齐全？
\end{enumerate}

\subsection{开放问题}

第一，能否让 high-resolution linear attention 保持 softmax attention 的细粒度关系？第二，模态专属参数与共享参数的最优比例是否随层深变化？第三，如何在 video 与 multi-image context 中动态分配 token budget？第四，in-context generation 需要怎样的数据格式与 task curriculum？第五，MoE、post-training 和 preference optimization 怎样与 diffusion backbone 协同？

这些问题说明未来进展不会只来自更大的 denoiser。更可能的突破是 backbone、数据、condition interface、系统 kernel、采样和 alignment 的共同设计。评价一个新 DiT 时，最有价值的问题不是“它用了哪个名字”，而是“它在什么约束下改变了哪条 Pareto frontier”。

\subsection{拓展阅读}

\begin{itemize}[itemsep=0.2em,topsep=0.2em]
  \item 官方 slides：\url{http://bit.ly/dit-cs25}
  \item 官方录像：\url{https://www.youtube.com/watch?v=vXtapCFctTI}
  \item 代表路线：DiT、PixArt-$\alpha$、SANA、Stable Diffusion 3/MMDiT、DiT-Air、ControlNet、LTX-Video、FuseDiT 与 OmniGen。
  \item 读论文时优先寻找 compute-matched ablation、端到端 latency、固定 prompt 评估、失败样本与完整 pipeline 成本。
\end{itemize}

\end{document}
